\chapter{Formal Appendix: Load-Bearing Mathematical Objects by Chapter}
\label{app:formal}

\begin{plainlanguage}
\textbf{In plain terms:} This appendix is a reference table, not a new argument. It lists, chapter by chapter, every formal object the book relies on, labeled honestly according to what kind of claim it actually is: a \emph{definition} (just naming something, no claim attached), a \emph{proposition} (an easy, near-immediate consequence of the definitions), a \emph{conjecture} (a precise, plausible, currently unproven claim), or a \emph{theorem} (something actually derived from stated assumptions, with a complete proof). Applying this discipline honestly turns out to matter: most of the book's formal content is definitions and conjectures, with only two genuine theorems in all eighteen chapters. That is the correct and healthy shape for a book whose subject is, in part, the difference between earned and borrowed authority --- a book that called everything a ``Theorem'' would be committing, in its own mathematics, exactly the error its own Chapter~5 exists to diagnose.
\end{plainlanguage}

Applied honestly, most of the book's formal content turns out to be definitions and conjectures, with a small number of genuine propositions and exactly one piece of mathematics that earns the word ``theorem'' without qualification --- joined by one more, a confirmation of corrected machinery in Chapter~\ref{ch:fidelity}, that is modest but also genuine. That distribution is the correct one for a book whose subject is, in part, the difference between earned and borrowed authority.

\section*{Part I --- The Problem of the Definition}

\subsection*{Chapter 1}

\textbf{Definition (Definition Space).} $\Theta = \{\theta_1, \theta_2, \ldots\}$, the space of candidate definitions of intelligence, equipped with a conceptual-distance metric $d_\Theta(\theta_i, \theta_j)$.

\textbf{Conjecture (Definition-Space Expansion).} In the historical record of a contested concept, $|\Theta_t|$ tends to increase before any convergence mechanism causes it to decrease. Motivates the need for elimination machinery in Part~II; does not establish that machinery is necessary.

\subsection*{Chapter 2}

\textbf{Definition (Hard/Soft Decomposition).} $\Theta_t = H_t \cup S_t$, where $H_t$ is the hard core and $S_t$ is the soft periphery.

\textbf{Conjecture (Differential Stability).} $dH_t/dt \to 0$ as evidence accumulates, while $dS_t/dt$ may remain nonzero indefinitely. The book's central organizing intuition, not yet a derived result --- it depends on what ``evidence accumulation'' formally does to $H$ and $S$, which is not specified until Chapter~\ref{ch:fidelity}.

\section*{Part II --- The Formal Machinery}

\subsection*{Chapter 3}

\textbf{Definition (Naive Elimination).} $\Theta_n = \bigcap_{i=1}^n C_i$, where each $C_i$ is a perspective-supplied constraint.

\textbf{Proposition (Monotonicity).} For pure intersection, $\Theta_{n+1} \subseteq \Theta_n$ for all $n$. Immediate from the definition of intersection.

\textbf{Definition (Perspective, added to Chapter 3).} A perspective $P$ is any source --- observer, culture, institution, developmental stage, research program, or system --- capable of producing evidence bearing on $\Theta$. Each perspective induces a constraint via $\Phi : \{\text{perspectives}\} \to \{\text{constraints on } \Theta\}$, with $C_i = \Phi(P_i)$. This closes a gap in an earlier draft: a book titled around \emph{perspectival convergence} had constraints ($C_i$) but no formal object named ``perspective'' until this definition was added.

\subsection*{Chapter 4}

\textbf{Definition (Belief Measure).} $\mu_t : \Theta \to [0,\infty)$, a plausibility assignment over candidate definitions, updated by $\mu_{t+1} = U(\mu_t, C_t)$ for some update rule $U$, rather than by set intersection.

\textbf{Theorem (Irreversible Elimination).} If a sequence of updates is strictly monotonic and shrinking --- $\Theta_{t+1} \subseteq \Theta_t$ for all $t$ --- then mistaken elimination of the true candidate $\theta^*$ can never be repaired by any subsequent update in the sequence. \emph{The strongest piece of mathematics in the book}: a real theorem, with explicit assumptions, a complete and non-trivial proof, and a conclusion that does actual argumentative work.

\subsection*{Chapter 5}

\textbf{Definition (Fidelity Functional).} $w_i = f(d_i, s_i, c_i, q_i)$, where $d_i$ is causal directness, $s_i$ is severity calibration, $c_i$ is closure, and $q_i$ is source independence. Normalized form: $w_i = d_i \cdot s_i \cdot c_i \cdot q_i$, with each factor scaled to $[0,1]$.

\textbf{Definition (Soft Elimination, corrected).} Originally proposed as $\mu_{t+1}(\theta) = \mu_t(\theta)(1 - w_t e_t(\theta))$, which does not preserve total probability mass under repeated application. Corrected via renormalization: $\mu_{t+1}(\theta) = \mu_t(\theta)(1 - w_t e_t(\theta)) / Z_t$, with $Z_t = \sum_\theta \mu_t(\theta)(1 - w_t e_t(\theta))$.

\textbf{Theorem (Fidelity Stability, under the corrected update).} If $0 \leq w_t e_t(\theta) \leq 1$ for all $\theta$ and $Z_t > 0$, then $\mu_t(\theta) \geq 0$ and $\sum_\theta \mu_t(\theta) = 1$ for all $t$. A real, if modest, theorem --- it confirms the corrected update rule behaves well, which the original uncorrected version did not.

\textbf{Conjecture (Misweighted Elimination).} $L = \mathbb{E}[(\hat\mu - \mu^*)^2]$; for correctly calibrated constraints, $\partial L/\partial w < 0$. The formal target the rest of the book gestures toward whenever it invokes ``low-fidelity counterfoils mistaken for hard eliminations,'' but not yet derived: requires a stated model relating $e_t(\theta)$ to the true $\mu^*$, which the book has not specified.

\section*{Part III --- How Convergence Actually Behaves}

\subsection*{Chapter 6}

\textbf{Definition (Interface Lag).} For coupled subsystems $S_i, S_j$ with closure times $t_i, t_j$, $\Delta t_{ij} = |t_i - t_j|$.

\textbf{Conjecture (Coupling--Lag Relation).} $\Delta t_{ij} \propto C_{ij}^{-1}$, where $C_{ij}$ measures coupling strength. Checked only qualitatively against one institutional case; no quantitative estimate of $C_{ij}$ attempted.

\textbf{Definition (Nominal vs.\ Operational Closure).} $N(t) \in \{0,1\}$, declared closure; $O(t) \in \{0,1\}$, actual operational closure. A \emph{closure illusion} is a state where $N(t)=1$ while $O(t)=0$.

\subsection*{Chapter 7}

\textbf{Definition (Subsystem-Indexed Repair).} Replace a single $R_n(X)$ with $R_n^{(k)}(X)$ for subsystem $k$.

\textbf{Principle (Anti-Vacuity).} A subsystem split into index $k$ is admissible only if there exists an independent measurement method $M_k$ for that subsystem, established by existing methodology rather than introduced solely to rescue a failed prediction. A methodological commitment, not a mathematical claim --- it has no proof because it isn't the kind of statement that has one.

\subsection*{Chapter 8}

\textbf{Definition (Cross-Domain Structure).} $\mathcal{S} = (\Delta t, N/O)$, the joint pattern of interface lag and closure status observed in a domain.

\textbf{Observation (Structural Resemblance, demoted from ``Cross-Domain Invariance'').} The same formal structure $\mathcal{S}$ was observed in two hand-selected domains. The originally proposed inequality $P(\mathcal{S}\mid D_1,D_2) > P(\mathcal{S}\mid D_1)P(\mathcal{S}\mid D_2)$ presupposes a sampling process that does not exist here; with $n=2$ non-randomly-chosen cases, no probability statement of this form is estimable.

\section*{Part IV --- The Pedagogy of Boundaries}

\subsection*{Chapter 9}

\textbf{Definition (Counterfoil Pair).} A pair $(a,c)$ with $a$ admissible, $c$ inadmissible, and distance $d(a,c)$ deliberately maximized subject to pedagogical relevance.

\textbf{Analogy (demoted from ``Margin Theorem'').} Counterfoil pedagogy resembles high-margin training contrasts in statistical classification. An evocative structural parallel, not a derived result --- the formal guarantees behind real margin-based learning theory have no established analogue for human concept acquisition via cultural pedagogy.

\subsection*{Chapter 10}

\textbf{Definition (Recognition and Implementation Competence).} $I_R = P(f(x)=y)$; $I_I = P(\pi(x)=a \mid f(x)=a)$. The conditional definition cleanly separates implementation failure from recognition failure by construction.

\textbf{Proposition (Dissociation).} High $I_R$ does not imply high $I_I$. Supported constructively, via the moral disengagement literature, not by formal proof.

\subsection*{Chapter 11}

\textbf{Definition (Curriculum Trajectory).} A stage-indexed map $s \mapsto (d_p(s), m(s))$.

\textbf{Conjecture (Interleaved Exposure, relabeled).} Effective exposure at stage $s$ is weighted: $E_s = \alpha_s I_R + \beta_s I_I$, with $\alpha_s > \beta_s$ early. Motivated by, but not established by, roughly a dozen illustrative, non-randomly-sampled cases.

\subsection*{Chapter 12}

\textbf{Definition (Transfer Matrix).} $T = \begin{pmatrix} T_{RR} & T_{RI} \\ T_{IR} & T_{II} \end{pmatrix}$.

\textbf{Proposition (Orthogonality Failure).} If $T_{IR} > 0$ or $T_{RI} > 0$, then $I_R$ and $I_I$ are not orthogonal. True essentially by definition.

\textbf{Proposition (Dissociability).} If there exist population subsets with $I_R \gg I_I$ and others with $I_I \gg I_R$, then the two capacities are non-identical. Also near-definitional.

\subsection*{Chapter 13}

\textbf{Definition (Recoverable Compound, replacing ``Recursive Compression Operator / Decompression Theorem'').} A competence label $\ell$ is recoverable-as-compound if there exist measurable $(x_1,\ldots,x_n)$, $n>1$, and an aggregation map $g : \mathbb{R}^n \to \mathbb{R}$ such that $\ell \approx g(x_1,\ldots,x_n)$, with the $x_i$ not mutually redundant. Replaces an earlier, mathematically unsound formulation ($C : \mathbb{R}^n \to \mathbb{R}$ paired with a claimed inverse $C^{-1}$): a many-to-one map has no general inverse.

\textbf{Conjecture (Recursive Compression).} Whenever a culture names a competence with a single word, closer inspection tends to reveal that the label is recoverable-as-compound. Currently supported by exactly one worked case, pending Chapter~\ref{ch:creativity}.

\section*{Part V --- A Second Pass}

\subsection*{Chapter 14}

\textbf{Conjecture (Generalization Test).} Creativity is recoverable-as-compound in the same sense. \emph{Outcome: partially confirmed.} A recognition/implementation-like split recurs, but creativity carries additional axes (domain expertise, social uptake) the intelligence case did not surface.

\textbf{Definition (Structural Isomorphism, as a test criterion).} A map $\varphi$ from $I$-components to $C$-components ``preserves decomposition structure'' if $\varphi(I_R)$ corresponds to a recognition-type component of creativity and $\varphi(I_I)$ to an implementation-type component. Offered as the chapter's evaluation criterion, not a proven isomorphism.

\subsection*{Chapter 15}

\textbf{Conjecture (Compression Efficiency, with explicit open cost model).} Communication cost favors compression --- $L(\text{profile}) \gg L(\text{single label})$ --- and a label persists whenever savings in description length exceed the predictive error introduced: $\Delta L > \Delta E$. Gestures at a real formal tradition (rate--distortion theory, minimum description length) but lacks the cost model that would let the inequality be evaluated. Left explicitly as the chapter's open formalization target.

\section*{Part VI --- Resolving the Fork}

\subsection*{Chapter 16}

\textbf{Definition (Borrowed Authority).} $B = 1 - q$, where $q$ is the source-independence factor from Chapter~\ref{ch:fidelity}'s fidelity functional.

\textbf{Definition (Counterfoil Distortion Index).} $D = (\text{depicted severity}) / (\text{actual severity})$.

\textbf{Proposition (checked for circularity, relabeled from ``Fidelity Theorem'').} World-forced eliminations are characterized by high $w$; rhetorical eliminations by high $B = 1-q$. Partly definitional, since $q$ is one factor in $w$; the substantive claim is that rhetorically-amplified counterfoils tend to score low on $q$ \emph{and} $d$, $s$, $c$ simultaneously --- joint co-occurrence, not the definitional relationship, is the chapter's real claim.

\subsection*{Chapter 17}

\textbf{Definition (Weak-Realist Decomposition, restated without overclaiming formal apparatus).} $\Theta = H \cup S$, where $H$ is the set of candidates $\theta \in \Theta$ that survive arbitrarily extended sequences of high-fidelity elimination. An earlier draft expressed $H$ as $\lim_{t\to\infty}$ of a set-valued process; set aside because such a limit requires specifying what converges and in what topology, which the book neither builds nor needs.

\textbf{Conjecture (Persistence, relabeled from ``Theorem'').} Elements of $H$ remain invariant under a (currently undefined) operation of ``cultural reparameterization.'' Cannot be a theorem because the reparameterization operation has not been formally defined.

\textbf{Definition (Survival vs.\ Corroboration; reordering principle for Section 17.4).} Survival: $\mu_t(\theta)$ stays bounded away from $0$ because nothing has eliminated it (purely negative --- the standing this book's Chapters 3--5 machinery already supports). Corroboration: $\mu_t(\theta)$ rises because independent perspectives actively recover it. Distinguished explicitly because the original draft of this section did not separate them, risking collapse of corroboration into mere consensus.

\textbf{Definition (Convergent Corroboration).} A candidate $\theta$ is convergently corroborated when recoverable across observational regimes --- independently re-derived by sufficiently independent perspectives, especially when separated by enough time or method that the later one could not have been shaped by the earlier one. Stated as the section's lead definition rather than left implicit until after the formalism.

\textbf{Definition (Independence, with two named sub-conditions).} $\mathrm{Ind}(i,j) \in [0,1]$, estimated from shared lineage history. Two conditions push it toward maximum: \emph{methodological distinctness} and \emph{temporal impossibility of contamination} (when the confirming methodology postdates the original claim, no shaping --- conscious or otherwise --- is possible; stronger than pre-registration). Flagged as the construct the entire section actually depends on: without a demanding independence criterion, corroboration collapses into consensus, fully vulnerable to the Chapter~16 borrowed-authority failure mode.

\textbf{Definition (Endorsement, Recovery Count, Pairwise/Aggregate Corroboration; Section 17.4).} $\varepsilon_i(\theta)$ measures whether a perspective specifically predicts $\theta$ rather than merely permitting it. $R(\theta,t)$ counts sufficiently independent perspective pairs that each specifically endorse $\theta$ above threshold --- the direct positive counterpart to an elimination count, distinct from and prior to the weighted $K_{ij}(\theta) = \varepsilon_i(\theta)\varepsilon_j(\theta)\mathrm{Ind}(i,j)$, aggregated as $K(\theta,t)$. Introduced to represent positive confirmation across independent perspectives, a relationship the purely negative elimination machinery of Chapters~3--5 cannot express. Illustrated by the Jakobson--Trubetzkoy distinctive-features case (theoretical claim, 1930s) and its independent confirmation by intracranial electrophysiology (2010s), as discussed by Yosef Grodzinsky. Computable in principle (e.g., via adjusted Rand index between independently-derived clusterings); not computed in this book.

\textbf{Definition (Belief-Field Lagrangian and Hamiltonian, proposed dynamical model --- explicitly secondary).} $\mathcal{L} = \frac{\kappa}{2}\dot\mu^2 - \frac{\sigma}{2}(\nabla_\theta\mu)^2 - V\mu$, with instability potential $V$ combining elimination pressure and corroboration; Legendre transform gives $\mathcal{H}$. This layer is explicitly framed as a formal description of a phenomenon the operational definitions above have already established, not as the source of the section's claim --- a reader unpersuaded by the dynamics loses none of the section's substance.

\textbf{Conjecture (Overdamped Reduction).} The Euler--Lagrange dynamics, in the slow-revision limit, plausibly reduce to something structurally close to Chapter~5's corrected soft-elimination update. Not rigorously confirmed; the discretization/renormalization needed to establish exact correspondence has not been carried out.

\subsection*{Chapter 18}

\textbf{Proposition (Reflexive Convergence, relabeled from ``Theorem'').} The sequence of models the book itself produces, $\Theta_0 \to \Theta_1 \to \Theta_2 \to \cdots$, instantiates the same elimination-and-repair structure defined formally in Chapters~3--5, as demonstrated concretely at Chapters~4, 7, 11, 12, and~14 (the last partially). A proposition supported by direct demonstration --- pointing at the specific instances where it happened --- rather than a theorem with independent assumptions and proof.

\section*{Summary table}

\begin{center}
\begin{tabular}{c c c c c}
\hline
\textbf{Chapter} & \textbf{Definitions} & \textbf{Propositions} & \textbf{Conjectures} & \textbf{Theorems} \\
\hline
1  & 1 & --- & 1 & --- \\
2  & 1 & --- & 1 & --- \\
3  & 2 & 1 & --- & --- \\
4  & 1 & --- & --- & \textbf{1 (genuine)} \\
5  & 2 & --- & 1 & \textbf{1 (genuine)} \\
6  & 2 & --- & 1 & --- \\
7  & 1 & --- & --- & --- (1 Principle) \\
8  & 1 & --- & --- & --- (1 Observation) \\
9  & 1 & --- & --- & --- (1 Analogy) \\
10 & 1 & 1 & --- & --- \\
11 & 1 & --- & 1 & --- \\
12 & 1 & 2 & --- & --- \\
13 & 1 & --- & 1 & --- \\
14 & 1 & --- & 1 & --- \\
15 & --- & --- & 1 & --- \\
16 & 2 & 1 & --- & --- \\
17 & 9 & --- & 2 & --- \\
18 & --- & 1 & --- & --- \\
\hline
\end{tabular}
\end{center}

Two genuine theorems out of eighteen chapters, both load-bearing (Chapter~\ref{ch:nonmonotonic} motivates the entire move to soft elimination; Chapter~\ref{ch:fidelity} confirms the corrected version of that machinery is well-behaved), surrounded by a much larger population of definitions and conjectures. That is the honest shape of the project at this stage, and it is a healthier one than a book with eighteen chapters each claiming a ``Theorem'' would have been --- a reader can now see exactly where the book has actually proven something, where it has only defined terms, and where it is making a precise, falsifiable bet it has not yet cashed in.
